% GENERATED FILE. Edit canonical lessons, then run npm run generate:books.
% canonical-source-hash: fnv1a64:b0ed2bd6e8ff5c6a
% canonical-lessons: TE-C233-kanduva, TE-C233-budaga, TE-C233-tuppu, TE-C233-masi, TE-C233-boggu

\chapter{Shoulder cloth, Bubble, Rust, Soot, Charcoal}
\label{ch:te-shoulder-cloth-bubble-rust-soot-charcoal}
\hlchaptermodality{233}

\begin{chapteropening}
\textbf{By the end of this chapter:} \emph{I can say a shoulder cloth, a bubble, rust, soot and charcoal.}

Complete the last lesson of chapter 233: I can say a shoulder cloth, a bubble, rust, soot and charcoal.
\end{chapteropening}

\section[\texorpdfstring{\te{కండువా} (kaṇḍuvā)}{కండువా (kaṇḍuvā)}]{\te{కండువా} (kaṇḍuvā) — a shoulder cloth, a scarf}

\label{lesson:TE-C233-kanduva}



\begin{quote}
Before the new one: say the Telugu for a barrel, then the Telugu for a mug.
\end{quote}

\subsection*{You'll want to know: \te{కండువా}}
\textbf{\te{కండువా}} — \emph{kaṇḍuvā} — \textquotedblleft{}a shoulder cloth, a scarf\textquotedblright{}.

\begin{grammarlens}[title={What you've built}]
One more word for this chapter.
\end{grammarlens}

\subsection*{Your turn}
\begin{itemize}
\raggedright
  \item \emph{Say it:} \emph{kaṇḍuvā}
  \item \emph{Say it:} \emph{kaṇḍuvā}, once more
  \item \emph{Say it:} say \emph{kaṇḍuvā}
  \item \emph{Recall:} say the Telugu for a barrel, then the Telugu for a mug, then say \emph{kaṇḍuvā} again
\end{itemize}

\begin{tcolorbox}[breakable,colback=teal!4,colframe=teal!35!black,title={Before you move on}]
What does \te{కండువా} mean? (\textquotedblleft{}A shoulder cloth, a scarf\textquotedblright{}.) Say it once more.
\end{tcolorbox}
\section[\texorpdfstring{\te{బుడగ} (buḍaga)}{బుడగ (buḍaga)}]{\te{బుడగ} (buḍaga) — a bubble}

\label{lesson:TE-C233-budaga}



\begin{quote}
Before the new one: say the Telugu for a mug, then the Telugu for a shoulder cloth.
\end{quote}

\subsection*{You'll want to know: \te{బుడగ}}
\textbf{\te{బుడగ}} — \emph{buḍaga} — \textquotedblleft{}a bubble\textquotedblright{}.

\begin{grammarlens}[title={What you've built}]
One more word for this chapter.
\end{grammarlens}

\subsection*{Your turn}
\begin{itemize}
\raggedright
  \item \emph{Say it:} \emph{buḍaga}
  \item \emph{Say it:} \emph{buḍaga}, once more
  \item \emph{Say it:} say \emph{buḍaga}
  \item \emph{Recall:} say the Telugu for a mug, then the Telugu for a shoulder cloth, then say \emph{buḍaga} again
\end{itemize}

\begin{tcolorbox}[breakable,colback=teal!4,colframe=teal!35!black,title={Before you move on}]
What does \te{బుడగ} mean? (\textquotedblleft{}A bubble\textquotedblright{}.) Say it once more.
\end{tcolorbox}
\section[\texorpdfstring{\te{తుప్పు} (tuppu)}{తుప్పు (tuppu)}]{\te{తుప్పు} (tuppu) — rust}

\label{lesson:TE-C233-tuppu}



\begin{quote}
Before the new one: say the Telugu for a shoulder cloth, then the Telugu for a bubble.
\end{quote}

\subsection*{You'll want to know: \te{తుప్పు}}
\textbf{\te{తుప్పు}} — \emph{tuppu} — \textquotedblleft{}rust\textquotedblright{}.

\begin{grammarlens}[title={What you've built}]
One more word for this chapter.
\end{grammarlens}

\subsection*{Your turn}
\begin{itemize}
\raggedright
  \item \emph{Say it:} \emph{tuppu}
  \item \emph{Say it:} \emph{tuppu}, once more
  \item \emph{Say it:} say \emph{tuppu}
  \item \emph{Recall:} say the Telugu for a shoulder cloth, then the Telugu for a bubble, then say \emph{tuppu} again
\end{itemize}

\begin{tcolorbox}[breakable,colback=teal!4,colframe=teal!35!black,title={Before you move on}]
What does \te{తుప్పు} mean? (\textquotedblleft{}Rust\textquotedblright{}.) Say it once more.
\end{tcolorbox}
\section[\texorpdfstring{\te{మసి} (masi)}{మసి (masi)}]{\te{మసి} (masi) — soot}

\label{lesson:TE-C233-masi}



\begin{quote}
Before the new one: say the Telugu for a bubble, then the Telugu for rust.
\end{quote}

\subsection*{You'll want to know: \te{మసి}}
\textbf{\te{మసి}} — \emph{masi} — \textquotedblleft{}soot\textquotedblright{}.

\begin{grammarlens}[title={What you've built}]
One more word for this chapter.
\end{grammarlens}

\subsection*{Your turn}
\begin{itemize}
\raggedright
  \item \emph{Say it:} \emph{masi}
  \item \emph{Say it:} \emph{masi}, once more
  \item \emph{Say it:} say \emph{masi}
  \item \emph{Recall:} say the Telugu for a bubble, then the Telugu for rust, then say \emph{masi} again
\end{itemize}

\begin{tcolorbox}[breakable,colback=teal!4,colframe=teal!35!black,title={Before you move on}]
What does \te{మసి} mean? (\textquotedblleft{}Soot\textquotedblright{}.) Say it once more.
\end{tcolorbox}
\section[\texorpdfstring{\te{బొగ్గు} (boggu)}{బొగ్గు (boggu)}]{\te{బొగ్గు} (boggu) — charcoal, coal}

\label{lesson:TE-C233-boggu}



\begin{quote}
Before the new one: say the Telugu for rust, then the Telugu for soot.
\end{quote}

\subsection*{You'll want to know: \te{బొగ్గు}}
\textbf{\te{బొగ్గు}} — \emph{boggu} — \textquotedblleft{}charcoal, coal\textquotedblright{}.

\begin{grammarlens}[title={What you've built}]
One more word for this chapter.
\end{grammarlens}

\subsection*{Your turn}
\begin{itemize}
\raggedright
  \item \emph{Say it:} \emph{boggu}
  \item \emph{Say it:} \emph{boggu}, once more
  \item \emph{Say it:} say \emph{boggu}
  \item \emph{Recall:} say the Telugu for rust, then the Telugu for soot, then say \emph{boggu} again
\end{itemize}

\begin{tcolorbox}[breakable,colback=teal!4,colframe=teal!35!black,title={Before you move on}]
What does \te{బొగ్గు} mean? (\textquotedblleft{}Charcoal, coal\textquotedblright{}.) Say it once more.
\end{tcolorbox}
